\section{Related work}
\label{sec:related}

\paragraph{Federated unlearning.}
FU methods fall into four families \cite{jeong2024sok,zhao2023fu-survey}.
\emph{History-based rollback and calibration} methods store per-round client updates
and reconstruct a model without the departing client: FedEraser calibrates stored
update directions \cite{liu2021federaser}, SIFU rolls back to a pre-influence
checkpoint with provable guarantees \cite{fraboni2024sifu}, and FedRecovery adds
calibrated DP noise to the recovered model \cite{zhang2023fedrecovery}. All require
$O(N R |\theta|)$ server-side storage. \emph{History-free optimisation} methods
(Halimi et al.\ \cite{halimi2022}, FedOSD \cite{fedosd2025}) perform constrained
gradient ascent on the departing client's data but require that client online at
request time. \emph{Weight manipulation and pruning} methods, i.e., weight negation
(NoT \cite{khalil2025not}), unlearning-during-learning (FedAU \cite{gu2024fedau}), and
class-discriminative pruning \cite{wang2022class-pruning}, avoid both storage and
cooperation but are heuristic. \emph{Exact-by-construction} methods such as
asynchronous cluster-based FU \cite{su2023async-fu} are the FL analogue of SISA
\cite{bourtoule2021sisa}. FORGE lies outside these four families: forgetting is a
property of the training dynamics, closed at request time by a non-interactive
subtraction on the head.

\paragraph{DP as unlearning and federated DP.}
Guo et al.\ \cite{guo2020certified} showed that an $(\varepsilon,\delta)$-DP model is
indistinguishable from one trained without a specific record. Sekhari et al.\
\cite{sekhari2021remember} separated deletion budget from privacy budget; Neel et al.\
\cite{neel2021descent} gave gradient-based deletion algorithms; Gupta et al.\
\cite{gupta2021adaptive} handled adversarial deletion sequences and extended the
analysis to non-convex losses. Ginart et al.\ \cite{ginart2019making} formalised the
deletion problem. All of these are centralised and record-level; forward forgetting
(Def.~\ref{def:forward-forgetting}) is the FL-native, client-level, time-indexed
analogue. In the federated setting, DP-FedAvg \cite{mcmahan2018learning} gives a
once-for-all client-level DP bound via clipping and server noise; FORGE inserts
temporal renewal via the decay schedule. Rényi-DP accounting
\cite{mironov2017renyi,abadi2016dpsgd} and the continual-release framework
\cite{chan2011private,dwork2010differential} are the direct antecedents of our
cumulative-noise analysis.

\paragraph{Centralised unlearning.}
Approximate unlearning via influence functions and Newton updates
\cite{izzo2021approximate,mehta2022deep}, representation-level forgetting
\cite{golatkar2020eternal,golatkar2021mixed,tarun2023fast}, class-level gradient-ascent
unlearning \cite{tarun2023fast,chundawat2023badteacher}, and data-free unlearning
\cite{chundawat2023zeroshot} inform the design vocabulary but assume a central dataset.
FORGE's head-only subtraction is closest to feature scrubbing
\cite{golatkar2020eternal}, adapted to use a client-committed DP summary rather than a
server-held estimate.

\paragraph{Memorisation and attack-channel leakage.}
Wei et al.\ \cite{wei2026fedmemprune} argue FU should remove memorised information
attributable to the forgotten data while preserving overlapping patterns; FORGE's
decay handles the shared-signal component and the summary channel handles the
memorised residual. On the attack side, Chen et al.\ \cite{chen2021unlearningleaks}
showed that comparing a model with its unlearned version enables an MIA that
outperforms the classical attack, and FUIA \cite{fuia2025} exploits this in the
federated setting. FORGE's protocol never publishes a before/after pair to any party
except the departing client themselves; the erasure procedure emits a hash of the
pre-committed summary plus a proof-of-execution transcript. Section~\ref{sec:results-fuia}
reports FUIA quality against FedRecovery and FORGE.

\paragraph{Federated distillation and privacy-preserving FU.}
Federated distillation methods \cite{li2019fedmd,he2020fedgkt,lin2020feddf} transmit
logits on a public probe set rather than gradients; our summary channel, when
instantiated with DP logits, is a privacy-noised variant of this transmission
repurposed as an unlearning substrate. Privacy-preserving FU protocols
\cite{starfish2024,fedadp2024,varshney2025plausible} protect the unlearning
\emph{computation} (two-party computation, adaptive DP, cluster-level plausible
deniability); FORGE is complementary in that it ensures the unlearning
\emph{substrate} is DP at the point of contribution, so the erasure step itself is
post-processing. On the verification side, Starfish \cite{starfish2024} runs the
unlearning computation under two-party computation with a fully adversarial
verifier; VeriFi \cite{gao2022verifi} assumes the verifier is a trusted third
party. FORGE sits between these two models: the authority is semi-honest (follows
the protocol, does not share information), which admits the hash-chain
verification described in Section~\ref{sec:method-erasure} without requiring MPC
machinery. For deployments that need a fully adversarial verifier, the hash
verification can be replaced by a succinct zero-knowledge proof of the same
relation.

\paragraph{Continual learning and unlearning.}
Continual learning surveys \cite{parisi2019cl-survey,delange2022cl-survey} treat
catastrophic forgetting as a failure mode; FCL work \cite{yoon2021,dong2022fcl}
extends this to federated settings. Continual \emph{unlearning} is thin
\cite{shibata2021learning-to-forget}. FORGE's decayed aggregation inherits the
temporal decay of continual learning and repurposes it as the substrate for a formal
erasure guarantee.
